% Options for packages loaded elsewhere
\PassOptionsToPackage{unicode}{hyperref}
\PassOptionsToPackage{hyphens}{url}
\PassOptionsToPackage{dvipsnames,svgnames,x11names}{xcolor}
\PassOptionsToPackage{space}{xeCJK}
%
\documentclass[
  11pt,
  a4paper,
]{article}
\usepackage{amsmath,amssymb}
\usepackage{iftex}
\ifPDFTeX
  \usepackage[T1]{fontenc}
  \usepackage[utf8]{inputenc}
  \usepackage{textcomp} % provide euro and other symbols
\else % if luatex or xetex
  \usepackage{unicode-math} % this also loads fontspec
  \defaultfontfeatures{Scale=MatchLowercase}
  \defaultfontfeatures[\rmfamily]{Ligatures=TeX,Scale=1}
\fi
\usepackage{lmodern}
\ifPDFTeX\else
  % xetex/luatex font selection
  \setmainfont[]{DejaVu Serif}
  \setmonofont[]{DejaVu Sans Mono}
  \ifXeTeX
    \usepackage{xeCJK}
    \setCJKmainfont[]{Noto Serif CJK SC}
          \fi
  \ifLuaTeX
    \usepackage[]{luatexja-fontspec}
    \setmainjfont[]{Noto Serif CJK SC}
  \fi
\fi
% Use upquote if available, for straight quotes in verbatim environments
\IfFileExists{upquote.sty}{\usepackage{upquote}}{}
\IfFileExists{microtype.sty}{% use microtype if available
  \usepackage[]{microtype}
  \UseMicrotypeSet[protrusion]{basicmath} % disable protrusion for tt fonts
}{}
\makeatletter
\@ifundefined{KOMAClassName}{% if non-KOMA class
  \IfFileExists{parskip.sty}{%
    \usepackage{parskip}
  }{% else
    \setlength{\parindent}{0pt}
    \setlength{\parskip}{6pt plus 2pt minus 1pt}}
}{% if KOMA class
  \KOMAoptions{parskip=half}}
\makeatother
\usepackage{xcolor}
\usepackage[margin=22mm]{geometry}
\usepackage{color}
\usepackage{fancyvrb}
\newcommand{\VerbBar}{|}
\newcommand{\VERB}{\Verb[commandchars=\\\{\}]}
\DefineVerbatimEnvironment{Highlighting}{Verbatim}{commandchars=\\\{\}}
% Add ',fontsize=\small' for more characters per line
\newenvironment{Shaded}{}{}
\newcommand{\AlertTok}[1]{\textcolor[rgb]{1.00,0.00,0.00}{\textbf{#1}}}
\newcommand{\AnnotationTok}[1]{\textcolor[rgb]{0.38,0.63,0.69}{\textbf{\textit{#1}}}}
\newcommand{\AttributeTok}[1]{\textcolor[rgb]{0.49,0.56,0.16}{#1}}
\newcommand{\BaseNTok}[1]{\textcolor[rgb]{0.25,0.63,0.44}{#1}}
\newcommand{\BuiltInTok}[1]{\textcolor[rgb]{0.00,0.50,0.00}{#1}}
\newcommand{\CharTok}[1]{\textcolor[rgb]{0.25,0.44,0.63}{#1}}
\newcommand{\CommentTok}[1]{\textcolor[rgb]{0.38,0.63,0.69}{\textit{#1}}}
\newcommand{\CommentVarTok}[1]{\textcolor[rgb]{0.38,0.63,0.69}{\textbf{\textit{#1}}}}
\newcommand{\ConstantTok}[1]{\textcolor[rgb]{0.53,0.00,0.00}{#1}}
\newcommand{\ControlFlowTok}[1]{\textcolor[rgb]{0.00,0.44,0.13}{\textbf{#1}}}
\newcommand{\DataTypeTok}[1]{\textcolor[rgb]{0.56,0.13,0.00}{#1}}
\newcommand{\DecValTok}[1]{\textcolor[rgb]{0.25,0.63,0.44}{#1}}
\newcommand{\DocumentationTok}[1]{\textcolor[rgb]{0.73,0.13,0.13}{\textit{#1}}}
\newcommand{\ErrorTok}[1]{\textcolor[rgb]{1.00,0.00,0.00}{\textbf{#1}}}
\newcommand{\ExtensionTok}[1]{#1}
\newcommand{\FloatTok}[1]{\textcolor[rgb]{0.25,0.63,0.44}{#1}}
\newcommand{\FunctionTok}[1]{\textcolor[rgb]{0.02,0.16,0.49}{#1}}
\newcommand{\ImportTok}[1]{\textcolor[rgb]{0.00,0.50,0.00}{\textbf{#1}}}
\newcommand{\InformationTok}[1]{\textcolor[rgb]{0.38,0.63,0.69}{\textbf{\textit{#1}}}}
\newcommand{\KeywordTok}[1]{\textcolor[rgb]{0.00,0.44,0.13}{\textbf{#1}}}
\newcommand{\NormalTok}[1]{#1}
\newcommand{\OperatorTok}[1]{\textcolor[rgb]{0.40,0.40,0.40}{#1}}
\newcommand{\OtherTok}[1]{\textcolor[rgb]{0.00,0.44,0.13}{#1}}
\newcommand{\PreprocessorTok}[1]{\textcolor[rgb]{0.74,0.48,0.00}{#1}}
\newcommand{\RegionMarkerTok}[1]{#1}
\newcommand{\SpecialCharTok}[1]{\textcolor[rgb]{0.25,0.44,0.63}{#1}}
\newcommand{\SpecialStringTok}[1]{\textcolor[rgb]{0.73,0.40,0.53}{#1}}
\newcommand{\StringTok}[1]{\textcolor[rgb]{0.25,0.44,0.63}{#1}}
\newcommand{\VariableTok}[1]{\textcolor[rgb]{0.10,0.09,0.49}{#1}}
\newcommand{\VerbatimStringTok}[1]{\textcolor[rgb]{0.25,0.44,0.63}{#1}}
\newcommand{\WarningTok}[1]{\textcolor[rgb]{0.38,0.63,0.69}{\textbf{\textit{#1}}}}
\usepackage{longtable,booktabs,array}
\usepackage{calc} % for calculating minipage widths
% Correct order of tables after \paragraph or \subparagraph
\usepackage{etoolbox}
\makeatletter
\patchcmd\longtable{\par}{\if@noskipsec\mbox{}\fi\par}{}{}
\makeatother
% Allow footnotes in longtable head/foot
\IfFileExists{footnotehyper.sty}{\usepackage{footnotehyper}}{\usepackage{footnote}}
\makesavenoteenv{longtable}
\setlength{\emergencystretch}{3em} % prevent overfull lines
\providecommand{\tightlist}{%
  \setlength{\itemsep}{0pt}\setlength{\parskip}{0pt}}
\setcounter{secnumdepth}{-\maxdimen} % remove section numbering
\ifLuaTeX
\usepackage[bidi=basic]{babel}
\else
\usepackage[bidi=default]{babel}
\fi
\babelprovide[main,import]{chinese}
\ifPDFTeX
\else
\babelfont{rm}[]{DejaVu Serif}
\fi
% get rid of language-specific shorthands (see #6817):
\let\LanguageShortHands\languageshorthands
\def\languageshorthands#1{}
\usepackage{fancyhdr}
\pagestyle{fancy}
\fancyhf{}
\fancyhead[L]{\small Order 15 / Third roots of unity}
\fancyhead[R]{\small Final supplement: 38 / 38}
\fancyfoot[C]{\thepage}
\setlength{\headheight}{15pt}
\setlength{\parindent}{0pt}
\setlength{\parskip}{0.45em}
\setcounter{tocdepth}{1}
\usepackage{fvextra}
\DefineVerbatimEnvironment{Highlighting}{Verbatim}{breaklines,commandchars=\\\{\},fontsize=\small}
\emergencystretch=2em
\ifLuaTeX
  \usepackage{selnolig}  % disable illegal ligatures
\fi
\usepackage{bookmark}
\IfFileExists{xurl.sty}{\usepackage{xurl}}{} % add URL line breaks if available
\urlstyle{same}
\hypersetup{
  pdftitle={15 阶三次单位根最大行列式：最后三层与全局上界},
  pdfauthor={研究续篇与可复演证书},
  pdflang={zh-CN},
  colorlinks=true,
  linkcolor={black},
  filecolor={Maroon},
  citecolor={Blue},
  urlcolor={blue},
  pdfcreator={LaTeX via pandoc}}

\title{15 阶三次单位根最大行列式：最后三层与全局上界}
\usepackage{etoolbox}
\makeatletter
\providecommand{\subtitle}[1]{% add subtitle to \maketitle
  \apptocmd{\@title}{\par {\large #1 \par}}{}{}
}
\makeatother
\subtitle{Q = 96, 99, 105 的完整排除 · 38 / 38 · 精确计算机辅助证明}
\author{研究续篇与可复演证书}
\date{2026 年 9 月 19 日}

\begin{document}
\maketitle

\subsection{1.
结论、范围与证明组织}\label{ux7ed3ux8bbaux8303ux56f4ux4e0eux8bc1ux660eux7ec4ux7ec7}

令 \(\omega^2+\omega+1=0\)，\(\mu_3=\{1,\omega,\omega^2\}\)，并设

\[G=HH^*,\qquad K=H^*H,\qquad Q=\sum_{i<j}|G_{ij}|^2,
\qquad H\in\mu_3^{15\times15}.\]

沿用原独立审计 {[}A17{]}、35 层续篇 {[}P35{]} 中的结构证明，以及已发表的
\(B_3(15,10)=12\) 码界 {[}TB20{]}，本续篇补齐
\(Q=96,99,105\)。合并这些结果得到

\[\boxed{\max_{H\in\mu_3^{15\times15}}|\det H|^2
=277868041444786176=2^{22}\cdot3^{20}\cdot19.}\]

等价地，最大行列式模为
\(120932352\sqrt{19}\)。原纪录矩阵达到这个值，故上界与下界相合。本结论确定最大值；没有声称已经分类所有等号矩阵，也没有声称最大化矩阵在通常等价关系下唯一。

这是带有精确有限证书的\textbf{计算机辅助数学证明}。本报告说明结构覆盖和各排除判据的数学依据；程序完成有限候选的生成、整数与有理运算、两种独立的全量列枚举，以及证书核验。它尚未被翻译为
Lean、Coq 等证明助手的形式化推导。本次没有重新执行 {[}TB20{]} 的码分类。

{[}P35{]} 的原文件、历史 35/38
状态表和原执行入口保留不变。最终状态由新的 \texttt{run\_final.py} 和
\texttt{results/final\_ledger.json}
汇总，避免把旧报告中的研究缺口直接改名成定理。

\subsection{2. 从 35
层到三个有限问题}\label{ux4ece-35-ux5c42ux5230ux4e09ux4e2aux6709ux9650ux95eeux9898}

{[}A17,P35{]} 已排除 \(Q\not\equiv0,6\pmod9\) 的情形，并证明所有
\(Q\ge171\) 的行列式平方小于纪录 \(B\)。在剩余 38 个允许壳层中，35
个已完成。因此只需排除满足 \(|\det H|^2>B\) 的 \(Q=96,99,105\) 矩阵。

严格反例的行、列颜色型各自只可能为 \((15)\)、\((14,1)\)、\((13,2)\) 或
\((13,1,1)\)。对于 \(Q=96,105\)，只可能出现 \((14,1)\) 或
\((13,2)\)；对于 \(Q=99\)，只可能出现 \((15)\) 或 \((13,1,1)\)。同色非零
Gram 项的平方模至少为 9；异色项平方模至少为 3。

把含 13 行的多数颜色主块记为
\(A\)，另两行称为枢纽。令其内部能量、交叉能量和枢纽之间内积的平方模分别为
\(e,c,h\)。先前的精确分量---Schur 界将三层的 size-13 余项限定为
\(e=0,9,18\)。\(e=0\) 会给出 13 个正交三次单位根行，由 {[}TB20{]}
排除；\(e=9\) 由 {[}P35, §9{]} 的单边缺陷引理排除。对于 \(e=18\)，只剩

\[ (e,c,h)=(18,78,Q-96). \]

两条内部边的平方模均为 9，故支撑只能为 \(P_3\) 加 10 个孤立点，或
\(2K_2\) 加 9
个孤立点。此次没有删去内部孤立点，也没有使用已经失效的行列颜色``同方向轨道''论证。

本续篇先排除所有 size-13 余项。随后，\(Q=96,105\)
的严格反例被迫在两侧都具有 \((14,1)\) 型，\(Q=99\)
被迫在两侧都为全同色。这一顺序保证了跨颜色配对的覆盖；不需要预先假定行
Gram 和列 Gram 处于同一颜色型。

\subsection{3.
必须满足的列方程与精确穷举}\label{ux5fc5ux987bux6ee1ux8db3ux7684ux5217ux65b9ux7a0bux4e0eux7cbeux786eux7a77ux4e3e}

\textbf{列方程引理。} 若可逆矩阵 \(H\) 满足 \(HH^*=G\)，则

\[H^*G^{-1}H=I.\]

因此每一列 \(v\) 都必须属于

\[\mathcal C(G)=\{v\in\mu_3^{15}:v^*G^{-1}v=1\}.\]

如果 \(\operatorname{span}_{\mathbb C}\mathcal C(G)\) 的维数小于
15，\(G\) 就不能由三次单位根矩阵分解。这个判据只利用一个 Gram
矩阵，不需要知道另一侧的颜色、支撑或相位。

列方程在整体乘三次单位根下不变，因此可以固定 \(v_1=1\)，完整检查
\(3^{14}=4,782,969\) 个向量。设 \(L>0\)
清除逆矩阵全部分母，\(J=LG^{-1}\)。程序使用 \(\mathbb Z[\omega]\)
中的整数系数，检查

\[v^*Jv=L.\]

第一种算法按三进制次序改变坐标，利用

\[q(v+d e_k)=q(v)+2\operatorname{Re}\!\left(\bar d(Jv)_k\right)+J_{kk}|d|^2\]

更新二次型，并同步更新 \(Jv\)。第二种算法为每个向量重新计算全部 105
个坐标对的二次型贡献，完全不使用增量更新。两种算法在本次全部 194 个 Gram
候选上，逐项得到相同的完整列列表。

每个候选还以任意精度 Eisenstein 运算检查
\(GJ=LI\)、行列式、能量和颜色同余。Numba 内核使用有符号 64 位整数；令
\(M\) 为 \(J\) 两个整数系数的最大绝对值，执行前检验
\(100\cdot15^2M<2^{63}\) 和
\(L<2^{63}\)，这是涵盖上述更新及直接求和中间量的保守界。

对被列空间排除的候选，程序额外提供一个明确的非零
\(w\in\mathbb Z[\omega]^{15}\)，并在完整允许列列表上检查
\(w^*v=0\)。因此最终证据包含 193
个可直接核验的湮灭向量，排除依据不限于一个浮点秩或未经展开的
\texttt{rank\ \textless{}\ 15} 返回值。

\subsection{4. size-13
的规范化与候选覆盖}\label{size-13-ux7684ux89c4ux8303ux5316ux4e0eux5019ux9009ux8986ux76d6}

记 \(\delta=1-\omega\)，行颜色为 \(r_i\)。同一个 Gram 内的正确同余是

\[G_{ij}\equiv(r_j-r_i)\delta\pmod{(3)}.\]

可整体共轭，使多数颜色与第一枢纽的颜色差为
1。由于所有交叉项平方模都是最小值
3，可以仅用三次单位根行缩放，将多数行到第一枢纽的项全部规范化为
\(\delta\)。第二组交叉项写为 \(\varepsilon\delta t_i\)，其中
\(t_i\in\mu_3\)，\(Q=99\) 时 \(\varepsilon=-1\)，另外两层
\(\varepsilon=1\)。再用第二枢纽的三次单位根缩放固定一个核心坐标的
\(t_i=1\)。

此处保持了 \(H\) 的三次单位根字母表。内部两条边分别枚举 \(3\mu_6\)
中的全部六个值，不把六次单位根行缩放误当成保留三次字母表的等价操作。对于
\(r\) 个内部孤立点，只需记录 \((n_0,n_1,n_2)\)，其和为
\(r\)；相同取值坐标间的排列属于允许的行置换。

第一枢纽到第二枢纽的项在 \(Q=96\) 时为零；\(Q=99\) 时枚举
\(\delta\mu_3\)；\(Q=105\) 时枚举全部平方模 9 的 Eisenstein
整数。核心边值、核心比值、孤立点计数和枢纽项均由有限循环生成。每个生成的
Gram 先接受精确 Schur 行列式计算，再接受 \(D>B\)、整数、\(3^{14}\mid D\)
及 Eisenstein 范数条件。

这些条件都是真实因子的必要条件。\(3^{14}\mid D\) 来自对 \(H\) 的 14
行减去第一行后提取 \(\delta^{14}\)；一个非负整数为 Eisenstein
范数，当且仅当每个 \(p\equiv2\pmod3\)
的有理素因子的指数均为偶数。生成程序直接使用这些判据，旧候选数值只用于一致性断言，未作为枚举输入来预先删去未知数值。

\begin{longtable}[]{@{}crrrr@{}}
\toprule\noalign{}
壳层 & \(P_3\) 规范代表 & \(2K_2\) 规范代表 & 合计 & 列空间排除 \\
\midrule\noalign{}
\endhead
\bottomrule\noalign{}
\endlastfoot
\(96\) & 20 & 44 & 64 & 64 \\
\(99\) & 20 & 47 & 67 & 67 \\
\(105\) & 8 & 3 & 11 & 10 \\
合计 & 48 & 94 & 142 & 141 \\
\end{longtable}

``规范代表''允许仍有等价冗余；表中数量并非不等价矩阵的分类数。所有 142
个候选都执行了两种全量列检查。唯一未被列空间维数排除的代表属于
\(Q=105\)，下一节给出其解析矛盾。

\subsection{\texorpdfstring{5. \(Q=105\)：模 4 范数矛盾与
\(16>15\)}{5. Q=105：模 4 范数矛盾与 16\textgreater15}}\label{q105ux6a21-4-ux8303ux6570ux77dbux76feux4e0e-1615}

最后一个 Gram 代表可写为

\[G_*=
\begin{pmatrix}A&z&z\\z^*&15&3\\z^*&3&15\end{pmatrix},\qquad
A=\operatorname{diag}(C,C,15I_9),\quad
C=\begin{pmatrix}15&3\\3&15\end{pmatrix},\quad z=\delta\mathbf1_{13}.\]

其行列式为
\(281238577200000000>B\)。它确实是一个必须进一步排除的抽象候选，单纯把数值与纪录比较无法结束证明。

设一个可能列向量为

\[v=(a,b,c,d,p_1,\ldots,p_9,u,v_0)^T\in\mu_3^{15},\qquad
s=a+b+c+d,\quad P=\sum_{j=1}^9p_j,\quad k=s/18+P/15.\]

由 Schur 补精确求逆得

\[v^*G_*^{-1}v
=x^*A^{-1}x+\frac{15}{392}|u+v_0-2\bar\delta k|^2
+\frac1{24}|u-v_0|^2,\]

其中 \(x\) 是前 13
个坐标。程序还以任意精度系数逐项核对这个二次型恒等式与完整逆矩阵一致。

假设同时有 \(a\ne b\) 和 \(c\ne d\)。每个不同符号对在 \(C^{-1}\)
上的贡献为 \(11/72\)，另外九个坐标贡献 \(9/15\)，从而

\[x^*A^{-1}x=2\cdot\frac{11}{72}+\frac9{15}=\frac{163}{180}.\]

如果 \(u\ne v_0\)，则 \(|u-v_0|^2=3\)，上式立即给出

\[v^*G_*^{-1}v\ge\frac{163}{180}+\frac3{24}
=\frac{371}{360}>1,\]

与列方程矛盾。如果 \(u=v_0\)，列方程则强制

\[\frac{15}{98}|u-\bar\delta k|^2=\frac{17}{180}.\]

清除分母并利用 \(\delta\bar\delta=3\)，得到

\[\boxed{N(30\delta u-5s-6P)=1666.}\]

左侧是 Eisenstein 整数的范数。对于 \(N(a+b\omega)=a^2-ab+b^2\)，模 4
的可能余数只有 \(0,1,3\)，而
\(1666\equiv2\pmod4\)。因此这一情形也不可能。两种枢纽情形都被排除，故每个可用列必须满足

\[a=b\quad\text{或}\quad c=d.\]

这条限制具有完整的代数证明，独立于列枚举。作为实现交叉检查，全部 5887
个规范化允许列中，两对同时不同的列数确为零：仅第一对不同为 2304
列，仅第二对不同为 2304 列，两对均相同为 1279 列。

现在回到整个矩阵。\(G_{*,12}=G_{*,34}=3\)。若一对长度 15
的三次单位根行在 \(a_0\) 个位置相同，则其实内积为
\(a_0-(15-a_0)/2\)。实内积为 3 强制 \(a_0=7\)，即每对行恰有 8
个不同位置。上面的列限制要求这两对行的不同位置互不相交；15
个位置却需要容纳 \(8+8=16\) 个位置，矛盾。

所以 \(G_*\) 无法分解，\(Q=105\) 的全部 size-13 分支排除。{[}P35{]}
已排除该层双 \((14,1)\) 有孤立点的严格反例；{[}A17{]}
的无孤立点分量界包含等号 \(B\)，并排除大于 \(B\) 的值。由此 \(Q=105\)
完整闭合。

\subsection{\texorpdfstring{6. \(Q=96\)：三个双 \((14,1)\)
支撑余项}{6. Q=96：三个双 (14,1) 支撑余项}}\label{q96ux4e09ux4e2aux53cc-141-ux652fux6491ux4f59ux9879}

size-13 已排除，因此两侧都必须是 \((14,1)\)。多数块的 14
个顶点若没有孤立点，至少需要七条内部边，总能量至少为
\(7\cdot9+42=105\)，超过当前的 96。因此只需考虑有内部孤立点的分支。

{[}P35, §10{]}
的双侧核支撑与平坦矩形秩约束，在完整小分量目录中只保留三个签名。新程序重新枚举具有这些签名的图支撑、范数与相位，确认其分别为平衡的
\(P_3\)、\(C_4\) 和
\(K_{2,3}\)，所有圈相位均可规范化为零。这里``平衡''只用于描述核心 Gram
边；后续原始三次单位根列测试保留全部实际边相位。

\begin{longtable}[]{@{}
  >{\raggedright\arraybackslash}p{(\columnwidth - 10\tabcolsep) * \real{0.1667}}
  >{\raggedleft\arraybackslash}p{(\columnwidth - 10\tabcolsep) * \real{0.1667}}
  >{\raggedleft\arraybackslash}p{(\columnwidth - 10\tabcolsep) * \real{0.1667}}
  >{\raggedleft\arraybackslash}p{(\columnwidth - 10\tabcolsep) * \real{0.1667}}
  >{\raggedleft\arraybackslash}p{(\columnwidth - 10\tabcolsep) * \real{0.1667}}
  >{\raggedleft\arraybackslash}p{(\columnwidth - 10\tabcolsep) * \real{0.1667}}@{}}
\toprule\noalign{}
\begin{minipage}[b]{\linewidth}\raggedright
多数块支撑
\end{minipage} & \begin{minipage}[b]{\linewidth}\raggedleft
孤立点数
\end{minipage} & \begin{minipage}[b]{\linewidth}\raggedleft
内部能量
\end{minipage} & \begin{minipage}[b]{\linewidth}\raggedleft
交叉能量
\end{minipage} & \begin{minipage}[b]{\linewidth}\raggedleft
规范相位---交叉状态
\end{minipage} & \begin{minipage}[b]{\linewidth}\raggedleft
\(D>B\) 范数候选
\end{minipage} \\
\midrule\noalign{}
\endhead
\bottomrule\noalign{}
\endlastfoot
\(P_3+C_4\) & 7 & 54 & 42 & 7776 & 52 \\
\(C_4\) & 10 & 36 & 60 & 4536 & 0 \\
\(K_{2,3}\) & 9 & 54 & 42 & 1296 & 0 \\
\end{longtable}

对于交叉能量 60 的情形，超过最小值 42 的 18
单位能量可置于核心或任意孤立点。程序保留一个平方模 21 的项、两个平方模
12
的项及各自全部余类轨道；孤立点仅按行置换合并，没有被删除。核心内所有六次单位相位都被作为
Gram 标签枚举。

52 个余下 Gram 全部执行两种列方程穷举，其允许列的张成维数均至多为
9，故不可能组成可逆的 15 列矩阵。每一个都有独立的整数 Eisenstein
湮灭向量。这完成 \(Q=96\) 的最后分支。

\subsection{\texorpdfstring{7.
\(Q=99\)：完整纯色目录与谱迹不等式}{7. Q=99：完整纯色目录与谱迹不等式}}\label{q99ux5b8cux6574ux7eafux8272ux76eeux5f55ux4e0eux8c31ux8ff9ux4e0dux7b49ux5f0f}

size-13 的 67 个候选已经全部排除，所以行、列 Gram 都必须是全同色。令

\[T=(G-15I)/3,\qquad S=(K-15I)/3,\qquad TH=HS.\]

\subsubsection{7.1
有孤立点和无孤立点的完整覆盖}\label{ux6709ux5b64ux7acbux70b9ux548cux65e0ux5b64ux7acbux70b9ux7684ux5b8cux6574ux8986ux76d6}

若 \(T\) 有 \(r\) 个支撑孤立点，对应的 \(r\) 行 \(H\) 是范数平方 15
的正交向量，且都属于 \(\ker S\)。因此

\[15(P_{\ker S})_{jj}\ge r\qquad\text{对每个坐标 }j.\]

支撑叶子会迫使其邻点在全部核向量中为零，与此不等式矛盾。当前支撑最多有
11 条边；一个没有孤立点也没有叶子的 15 顶点图至少有 15
条边。因此两侧要么都无孤立点，要么都含孤立点，后一种情形的非平凡分量都无叶子。

复用 {[}P35, §12{]} 中已经完整生成的 12 能量单位目录，再限制总能量为
11。七个顶点以内的连通支撑使用完整图册；较大支撑由全部非同构树加全部允许额外边生成，随后精确同构去重。每个目标连通图都有生成树，因此这一构造不会遗漏目标支撑。边范数、单位轨道与所有弦相位均完整枚举。

有孤立点的分支得到 435 个分量组合，行列式和范数筛后保留 71
个压缩的整图谱状态。相同谱合并时取最大的最小核投影对角，始终向保留候选的方向放宽。再要求两侧同谱并同时满足上述投影不等式，候选全部消失。

无孤立点的分支得到 632 个分量组合。超过 \(B\) 且满足范数条件的候选有 12
个整图谱。谱迹测试需要知道各连通分量的谱，故第二次组装保留全部分量分解，最终得到
14 个分量---谱组合和 16 对无序的同谱行列配对。仅保留 12
个整图谱中的任意代表不够；本次对全部 14 个分解及全部 16 个配对逐一处理。

\subsubsection{7.2
一条无需近似特征值的精确矩形上界}\label{ux4e00ux6761ux65e0ux9700ux8fd1ux4f3cux7279ux5f81ux503cux7684ux7cbeux786eux77e9ux5f62ux4e0aux754c}

\textbf{谱迹引理。} 设 \(I,J\) 分别是 \(T,S\)
中若干完整连通分量的顶点并集，\(X=H_{I,J}\)。记

\[p(x)=\gcd\bigl(\chi_{T_I}(x),\chi_{S_J}(x)\bigr)
=x^d+c_1x^{d-1}+\cdots.\]

若 \(d=0\)，约定 \(c_1=0\)。则必要条件为

\[\boxed{|I||J|=\|X\|_F^2\le15d-3c_1.}\]

\textbf{证明。} 连通分量封闭性给出 \(T_I X=X S_J\)。在两个 Hermitian
矩阵的正交特征分解下，\(X\)
只能在相同特征值的空间之间非零。对于公共特征值
\(\lambda\)，该块秩至多为两侧重数的较小者。此外
\(XX^*\preceq G_I=15I+3T_I\)，故该块的每个非零奇异值平方至多为
\(15+3\lambda\)。对全部公共特征值按最小重数求和，得到
\(15d+3\sum\lambda=15d-3c_1\)。另一方面，每个 \(X_{ij}\) 模长都为
1，故左端是矩形面积。证毕。

这一不等式对任意单位模因子已经成立，不限于三次单位根。它保留不同特征值各自的贡献，比``公共谱维数乘一个统一最大特征值''更强，且右端完全由整数多项式的最大公因子给出。

程序枚举每个配对中两侧所有非空分量并集，为全部 16
对都找到违反该不等式的矩形。此前较弱上界留下的三对，新的整数矛盾分别为

\[40\le30,\qquad20\le18,\qquad36\le30.\]

对应的公共多项式分别为 \(x^2-1\)、\(x-1\)、\(x^2-1\)。全部 16
个证书均记录左右分量编号、两侧特征多项式、最大公因子、矩形面积与谱迹上界。由此完成
\(Q=99\) 的全同色分支和整个壳层。

\subsection{8.
全局结论的合并与等号矩阵}\label{ux5168ux5c40ux7ed3ux8bbaux7684ux5408ux5e76ux4e0eux7b49ux53f7ux77e9ux9635}

三个新排除都针对 \(D>B\) 的严格反例。size-13
不可实现性先消除了不同颜色之间的可能配对；随后分别处理强制的双
\((14,1)\)
或双全同色情形。所有依赖的有孤立点分支都保留了其真实交叉约束。

最终聚合器检查三个余项与旧精确分量表之间的覆盖关系，核对 142 个 size-13
代表、52 个额外 \(Q=96\) 代表、193 个湮灭向量、\(Q=105\) 解析矛盾以及
\(Q=99\) 全部 16 对谱证书。与历史 35 个壳层和 \(Q\ge171\)
的无限尾部联合，覆盖全部可能的 \(H\)。

已知矩阵的精确行列式为

\[\det H_0=604661760+241864704\omega,\]

其范数就是 \(B\)。原审计使用 Eisenstein
Bareiss、有理高斯消元，以及直接计算 \(\det(H_0H_0^*)\)
三条途径核对，并确认行列能量均为 105。该矩阵的 225 个指数保留在
\texttt{order15\_mu3\_unified\_audit/data/benchmark.json}；附录 A
也给出其逐行编码。

所以最大值等于 \(B\)，其相对 Hadamard 界的行列式模比例约为
\(0.7965900126\)。这个比例来自已验证的矩阵与上界，不是随机搜索的经验估计。

\subsection{9.
完整复演、独立验证与工程边界}\label{ux5b8cux6574ux590dux6f14ux72ecux7acbux9a8cux8bc1ux4e0eux5de5ux7a0bux8fb9ux754c}

在交付包根目录安装 \texttt{requirements.txt}
中的四个固定版本依赖后运行：

\begin{Shaded}
\begin{Highlighting}[]
\ExtensionTok{python}\NormalTok{ run\_final.py}
\end{Highlighting}
\end{Shaded}

程序默认重新执行原独立审计、前一阶段全部目录生成及证明程序，再从头生成三个最后壳层的全部新证书。\texttt{-\/-cached-baseline}
只复用四个哈希核验过的前阶段目录；新候选、两种全量列枚举和原独立审计仍会运行。不得使用
\texttt{python\ -O}。

本次完整执行耗时及逐程序耗时记录在
\texttt{results/final\_replay\_summary.json}。新的总复演包括 13
个顶层步骤，全部退出码为零；旧 35 层在其中又完成一次全量复演。总耗时为
595.64 秒，此数值仅描述本次容器环境，不作为其他机器的运行时间承诺。

新的每种列算法均检查 \(194\times3^{14}=927,895,986\)
个向量赋值，两种算法合计 \(1,855,791,972\) 次。对于每个
Gram，它们输出的完整允许列序列完全一致。此外，194
个候选都接受任意精度行列式、颜色同余与逆矩阵恒等式核对，193
个湮灭向量在其对应全部允许列上逐项验证。

\texttt{provenance/final\_source\_sha256.json}
固定本次新增源文件；原来源清单仍单独检查。新的程序在缺文件、源哈希不符、候选覆盖不符、任意证明步骤失败或仍有余项时均报错。运行开始会把新的执行摘要置为未完成状态，因此中断不能留下一个看似属于本轮的成功摘要。原
35 层入口继续诚实地报告其历史范围。

检查结论的最短路径是先读本报告第 2、5、7 节，再核对
\texttt{last\_three/final\_certificate.json}
与完整日志。核验有限计算时应重新运行程序；单独查看 JSON
中的``完成''字段不能替代生成算法和数学覆盖说明。

\subsection{10.
外部依赖、引文勘误与未声称事项}\label{ux5916ux90e8ux4f9dux8d56ux5f15ux6587ux52d8ux8befux4e0eux672aux58f0ux79f0ux4e8bux9879}

{[}TB20{]} 的 Proposition 14 和 Table 1 给出 \(B_3(15,10)=12\)。若有 13
个彼此正交的三次单位根行，将符号映射为三元码字，每对码字的 Hamming
距离为 10，会违反此界。此用途只需已发表的上界，无需重新构造其最优码。

前一份 35 层报告误写了 {[}TB20{]} 的期号和页码。正确书目信息为
\emph{Journal of Mathematical and Computational Science} \textbf{10
(2020), No.~6, 2713--2721}，DOI
10.28919/jmcs/4964。该处引文定位现已通过出版商 PDF 核对；所用
\(B_3(15,10)=12\)
数值没有变化。旧报告作为历史文件保留，勘误在此明确记录。

本续篇不使用旧颜色轨道方向错误所支持的孤立点排除，也不把局部候选的不可实现性自动推广到未枚举的支撑。全局范围来自第
2 节的旧约化、此次完整候选覆盖和第 8
节的合并。关于所有等号矩阵的分类、等价类数量以及形式化证明助手验证，均未包含在本次结论中。

\subsection{参考文献与内部证书}\label{ux53c2ux8003ux6587ux732eux4e0eux5185ux90e8ux8bc1ux4e66}

{[}TB20{]} Todor Todorov and Galina Bogdanova. \textbf{Ternary
equidistant codes of length \(11\le n\le15\)}. \emph{Journal of
Mathematical and Computational Science} 10 (2020), No.~6,
pp.~2713--2721. DOI: 10.28919/jmcs/4964. 使用 Proposition 14 和 Table
1；未复演其外部码分类。

{[}A17{]}
本项目原独立审计：\texttt{order15\_mu3\_unified\_audit/}。提供纪录矩阵的三重精确验证、颜色及能量约化、无限尾部、最初
19 个壳层与无孤立点分量界。

{[}P35{]} 本项目《15 阶三次单位根最大行列式：新增 16
个壳层的证明》，2026-09-18，\texttt{report/PROOF\_zh.md} 及对应
PDF。提供单边缺陷引理、双侧核投影与秩约束、扩展纯色目录，以及前一阶段的
35/38 完整证明范围。

{[}F38{]} 本续篇：\texttt{last\_three/}
的精确程序与结果，\texttt{run\_final.py}
的完整复演记录，\texttt{results/final\_ledger.json} 的全部 38
层状态与最终定理。

\subsection{附录 A.
达到上界的矩阵}\label{ux9644ux5f55-a.-ux8fbeux5230ux4e0aux754cux7684ux77e9ux9635}

下面每一行是一个长度 15 的指数串，数字 \(e_{ij}\in\{0,1,2\}\) 对应矩阵项
\(H_{ij}=\omega^{e_{ij}}\)。这份数据来自原审计的纪录矩阵，没有另行搜索或替换：

\begin{Shaded}
\begin{Highlighting}[]
\NormalTok{001000102102020}
\NormalTok{200221220112021}
\NormalTok{201110011122201}
\NormalTok{022200010001001}
\NormalTok{102122000102202}
\NormalTok{020211201100200}
\NormalTok{010021110121200}
\NormalTok{012112021111020}
\NormalTok{201212120201100}
\NormalTok{000110200221012}
\NormalTok{000122111000121}
\NormalTok{210101002101111}
\NormalTok{111212100020011}
\NormalTok{210012212002002}
\NormalTok{110010020200221}
\end{Highlighting}
\end{Shaded}

把每行转为 15 个指数后即可重建 \(H_0\)。完整数据、原仓库来源与提交定位见
\texttt{data/benchmark.json}。本报告只需要一个达到上界的实例；其他等号实例是否属于同一等价类不影响最大值定理。

\end{document}
